\chapter{Biological Plant Modeling}
\label{ch:biological_plant_modeling}

% Core Research Question:
% How do we construct a rigorous, whole-brain biophysical ground-truth simulation
% testbed of focal epilepsy with multi-channel scalp sensing and actuation?

% Contents:
% - Neural Mass Dynamics: Three-population Jansen-Rit cortical column equations and synaptic filtering kernels.
% - Bifurcation Regimes: Parameterization of healthy background tissue, autonomous Epileptogenic Zone (EZ),
%   and vulnerable Propagation Zone (PZ).
% - Whole-Brain Network: 76-region structural connectome coupling with distance-dependent axonal conduction delays
%   (stochastic delay differential equation system).
% - Electrical Sensing Interface: Regional Local Field Potentials (LFPs) projected to 62 scalp EEG channels via
%   volume-conduction Leadfield matrix.
% - Electrical Actuation Interface: Transcranial scalp currents mapped to regional electric fields and somatic
%   Stimulation Drives via ROAST finite-element head modeling (TP9-CP5-Ex8 montage).
% - Signal Conditioning: Causal fourth-order Butterworth low-pass filtering (25 Hz cutoff) and decimation to 50 Hz discrete grid.
% - Plant Calibration: Empirical recalibration of global coupling factor K and noise intensity sigma for realistic
%   seizure escape and recruitment dynamics.

% Mapping from Previous Structure:
% - Merges the biophysical modeling theory from Chapter 3 (03_foundations.tex, Sec. 3.1:
%   Jansen-Rit equations, bifurcation regimes, connectome coupling, delays, and dipole/leadfield electrophysiology)
%   with the concrete simulation setup from Chapter 4 (04_methodology.tex, Sec. 4.1:
%   76-region TVB setup, EZ/PZ temporal-lobe partition, ROAST TP9-CP5-Ex8 montage, Heun integration,
%   and causal anti-aliasing filtering/decimation).
